%% LyX 2.4.1 created this file.  For more info, see https://www.lyx.org/.
%% Do not edit unless you really know what you are doing.
\documentclass[10pt,english,hebrew]{article}
\usepackage{amsmath}
\usepackage{fontspec}
\setmainfont[Ligatures=TeX]{David}
\setsansfont[Ligatures=TeX]{David}
\setmonofont{Ebrima}
\usepackage[a4paper]{geometry}
\geometry{verbose,tmargin=0.5cm,bmargin=0.5cm,lmargin=0.5cm,rmargin=0.5cm,headheight=0cm,headsep=0cm,footskip=2cm}
\setcounter{secnumdepth}{-2}
\usepackage[skip=0bp]{parskip}
\usepackage{units}
\usepackage{mathtools}

\makeatletter

%%%%%%%%%%%%%%%%%%%%%%%%%%%%%% LyX specific LaTeX commands.
%% Because html converters don't know tabularnewline
\providecommand{\tabularnewline}{\\}

\@ifundefined{date}{}{\date{}}
%%%%%%%%%%%%%%%%%%%%%%%%%%%%%% User specified LaTeX commands.
\usepackage{multicol}
\usepackage{titlesec}

\renewcommand{\normalsize}{\footnotesize} % Makes all text smaller

\titleformat{\section}[runin]{\large\bfseries}{\thesection}{1em}{}
\titleformat{\subsection}[runin]{\large\bfseries}{\thesubsection}{1em}{}
\titleformat{\subsubsection}[runin]{\small\bfseries}{\thesubsubsection}{1em}{}

\makeatother

\usepackage{polyglossia}
\setdefaultlanguage{hebrew}
\setotherlanguage{english}
\begin{document}
\begin{multicols}{3}
\setlength{\columnseprule}{0.5pt}
\setlength{\abovedisplayskip}{0.1pt}
\setlength{\belowdisplayskip}{0.1pt}
\raggedcolumns
\setlength{\parskip}{0pt}
\tolerance=10000

\section{נוסחאון מכניקה אנליטית {\LR{2025}}}

\subsection{מתמטיקה}

\textbf{טריגו}

$c^{2}=a^{2}+b^{2}-2ab\cos\gamma$\\
$\cos\left(\alpha\pm\beta\right)=\cos\alpha\cos\beta\mp\sin\alpha\sin\beta$\\
$\sin\left(\alpha\pm\beta\right)=\cos\alpha\sin\beta\pm\sin\alpha\cos\beta$\\
$\sin\alpha\pm\sin\beta=2\sin\left(\frac{\alpha\pm\beta}{2}\right)\cos\left(\frac{\alpha\mp\beta}{2}\right)$\\
ׂ$\cos\alpha+\cos\beta=2\cos\left(\frac{\alpha+\beta}{2}\right)\cos\left(\frac{\alpha-\beta}{2}\right)$\\
$\cos\alpha-\cos\beta=-2\sin\left(\frac{\alpha+\beta}{2}\right)\left(\frac{\alpha-\beta}{2}\right)$

\paragraph{אוסילטור הרמוני מרוסן ומאולץ:}

\[
\ddot{x}+2\gamma\dot{x}+\omega_{0}^{2}x=f\left(x\right)
\]
מתקבל הפתרון הפרטי:
\[
x_{p}=\frac{1}{\sqrt{2\pi}}\int_{-\infty}^{\infty}\frac{\hat{f}\left(\omega\right)e^{i\omega t}}{\omega_{0}^{2}-\omega^{2}+2i\gamma\omega}d\omega
\]
ועבור פונקציה מאלצת $f\left(t\right)=f_{0}\cos\left(\omega_{d}t+\phi\right)$:
\[
x_{p}=\frac{f_{0}\cos\left(\omega_{d}t+\phi\right)}{\sqrt{\left(\omega_{0}-\omega_{d}\right)^{2}+\left(2\gamma\omega_{d}\right)^{2}}}
\]
\[
\phi=\tan^{-1}\left(\frac{2\gamma\omega_{d}}{\omega_{0}^{2}-\omega_{d}^{2}}\right)
\]


\paragraph{טריק למציאת מטריצה הופכית:}

\[
A=\begin{pmatrix}a & b\\
c & d
\end{pmatrix}\Rightarrow A^{-1}=\frac{1}{\det A}\begin{pmatrix}d & -b\\
-c & a
\end{pmatrix}
\]
\textbf{סימן לוי צ'יויטה} $\epsilon_{ijk}$ מחליף סימן בכל פרמוטציה,
ומתאפס אם המספרים הם לא פרמוטציה. לדוג':
\[
\epsilon_{123}=1,\ \epsilon_{213}=-1,\ \epsilon_{231}=1,\ \epsilon_{221}=0
\]
\textbf{מכפלה וקטורית מוכללת}\ \ \ יכולה להכתב בצורה:
\[
\left(\vec{v}\times\vec{u}\right)_{k}=\epsilon_{ijk}v_{i}u_{j}
\]
\textbf{טרנספורם לג'נדר}\ \ \ נניח שיש לנו פונקציה הפיכה $f\left(x\right)$.
נגדיר:
\[
\frac{df}{dx}=f^{\prime}\left(x\right)\equiv s\left(x\right)
\]
ומכאן נקבל פונקציה חדשה:
\[
g\left(s\right)=-x\left(s\right)\ast s+f\left(x\left(s\right)\right)
\]
\textbf{הטרנספורם ההפוך}\ \ \ נניח שקיים $g\left(s\right)$ שהוא
טרנספורם של $f\left(x\right)$. נמצא את $x\left(s\right)$ המקיים:$\frac{\partial g}{\partial s}=-x\left(s\right)$.
נקבל את $f\left(x\right)$: $V\left(x\right)=g\left(s\left(x\right)\right)+xu\left(x\right)$\\
\textbf{פונקציה הומוגנית מדרגה-}\textenglish[variant=american]{\textbf{k}}\ \ \ היא
פונקציה המקיימת $f\left(a\vec{x}\right)=a^{k}f\left(\vec{x}\right)$
לכל $a$\\
\textbf{משפט אוילר לפונקציה הומוגנית}\ \ \ $f$ הומוגנית חיובית
מסדר $k$ אם“ם: $\vec{x}\cdot\nabla f\left(\vec{x}\right)=kf\left(\vec{x}\right)$
או בגרסה החד מימדית: $xf\left(x\right)=k\frac{\partial f}{\partial x}\left(x\right)$.\\
\textbf{קור אליפטיות}\ \ \ $x=ar\cos\theta,\ y=br\sin\theta$.\\
$dA=abrdrd\theta$. אינטגל על השטח: $\int_{0}^{2\pi}\int_{0}^{1}dA$\\
\textbf{מטריצת סיבוב}\ \ \ לדוגמה, מטריצת סיבוב $\phi$ רדיאנים
סביב ציר ה$z$ 
\[
R_{\phi}^{z}=\begin{pmatrix}\cos\phi & \sin\phi & 0\\
-\sin\phi & \cos\phi & 0\\
0 & 0 & 1
\end{pmatrix}
\]


\paragraph{זוויות אוילר}

הן דרך לתאר מצב של גוף קשיח במרחב. הסיבוב מתואר ע“י {\LR{3}} טרנספורמציות
סיבוב, עבור $\phi$, $\theta$, $\psi$:
\[
D=R_{\psi}^{z}R_{\theta}^{x}R_{\phi}^{z}
\]
\textbf{מהירות זוויתית בזוויות אוילר:}

$\vec{\Omega}=\begin{pmatrix}\omega_{1}\\
\omega_{2}\\
\omega_{3}
\end{pmatrix}=\begin{pmatrix}\dot{\phi}\sin\theta\sin\psi+\dot{\theta}\cos\psi\\
\dot{\phi}\sin\theta\cos\psi+\dot{\theta}\sin\psi\\
\dot{\psi}+\dot{\phi}\cos\theta
\end{pmatrix}$

\paragraph{סוגרי פואסון}

\[
\left\{ f,g\right\} _{pq}\coloneqq\sum\left(\frac{\partial f}{\partial p_{i}}\frac{\partial g}{\partial q_{i}}-\frac{\partial f}{\partial q_{i}}\frac{\partial g}{\partial p_{i}}\right)
\]
\[
\left\{ f,g\right\} =-\left\{ f,g\right\} ,\ \left\{ af+bg,h\right\} =a\left\{ f,h\right\} +b\left\{ g,h\right\} 
\]
\[
\left\{ q_{i},q_{j}\right\} =\left\{ p_{i},p_{j}\right\} =0,\ \ \left\{ p_{i},q_{j}\right\} =\delta_{ij}
\]
\[
\left\{ ab,c\right\} =a\left\{ b,c\right\} +b\left\{ a,c\right\} 
\]
\[
\left\{ f,\left\{ g,h\right\} \right\} +\left\{ g,\left\{ h,f\right\} \right\} +\left\{ h,\left\{ f,g\right\} \right\} =0
\]
\[
\left\{ f\left(u\right),g\left(v\right)\right\} =\left\{ u,v\right\} \frac{\partial f}{\partial u}\frac{\partial g}{\partial v}
\]
\[
\left\{ f\left(\vec{p}\right),g\left(\vec{q}\right)\right\} =\sum\frac{\partial f}{\partial p_{i}}\frac{\partial g}{\partial q_{i}}
\]
\[
\left\{ f\left(\vec{q},\vec{p}\right),q_{i}\right\} =\frac{\partial f}{\partial p_{i}},\ \ \left\{ f\left(\vec{q},\vec{p}\right),p_{i}\right\} =-\frac{\partial f}{\partial q_{i}}
\]
\[
\left\{ L_{i},L_{j}\right\} =-\sum_{k}\epsilon_{ijk}L_{k}
\]
\textbf{מציאת מינימום מקומי עבור $f\left(x,y\right)$:}\ \ \ \ נמצא
נקודה $\nabla f=\left(0,0\right)$. נחשב: $D=f_{xx}f_{yy}-f_{xy}^{2}$.
נסווג:\\
$\boldsymbol{D>0}$\textbf{:} $f_{xx}>0$ ---\textless{} מינימום.
$f_{xx}<0$---\textless{} מקסימום\\
$\boldsymbol{D<0}$\textbf{: }אוכף.

\subsection{מכניקה לא אנליטית}

\textbf{תנ“ז }$\vec{L}=\vec{r}\times\vec{p}$

\textbf{מומנט כוח} $\vec{\tau}=\vec{r}\times\vec{F}$

\textbf{מתקף ותנע )קווי וזוויתי(} $\vec{J}=\int_{t_{1}}^{t_{2}}\vec{F}dt=\Delta\vec{p}$
\\
עבור תנ“ז, נגדיר $\vec{J}=\int_{t_{1}}^{t_{2}}\tau dt=\Delta L$

\subsection{משוואות לגראנז'}

בפרק זה נעסוק הרבה בדרגות חופש. במערכת עם $N$ חלקיקים יש $3N$ דרגות
חופש, כאשר רק \textbf{אילוצים הולונומיים} מורידים את דרגות החופש.
נקבל שבמערכת יש $n=3N-K$ דרגות חופש בלתי תלויות, כאשר $K$ מספר האילוים
ההולונומיים.

\paragraph{אילוצים הולונומיים:}

אם קיימת משוואה אלגברית מהצורה
\[
F\left(\vec{r_{1}},\vec{r_{2}},...,\vec{r_{n}},t\right)
\]


\paragraph{הלגרנג'יאן:}

$\mathcal{L}=K-V$ \textbf{משוואת א.ל}:
\[
\frac{d}{dt}\left(\frac{\partial\mathcal{L}}{\partial\dot{q}_{i}}\right)-\frac{\partial\mathcal{L}}{\partial q_{i}}=0\ \ i\in\left\{ 1,...,n\right\} 
\]
אפשר להוסיף ללגרנג'יאן פונקציה שהיא נגזרת שלמה בזמן של פונקציה כללית
$F\left(q,t\right)$.

\paragraph{משוואת א.ל עבור ”מכניקה מכוללת“}

כלומר כאשר יש לגרנג'יאן: $\mathcal{L}=\mathcal{L}\left(q,\dot{q},\ddot{q},t\right)$
נקבל א.ל:
\[
\frac{d^{2}}{dt^{2}}\left(\frac{\partial\mathcal{L}}{\partial\ddot{q}}\right)\boxed{-}\frac{d}{dt}\left(\frac{\partial\mathcal{L}}{\partial\dot{q}}\right)\boxed{+}\frac{\partial L}{\partial q}=0
\]


\paragraph{אנרגיה מוכללת}

$E=\sum_{i}\dot{q_{i}}\frac{\partial\mathcal{L}}{\partial\dot{q_{i}}}-L$\\
אם אין שינוי בלגרנג'יאן לאורך זמן, אז קיים שימור אנרגיה במערכת:$\frac{\partial L}{\partial t}=0\Rightarrow\frac{dE}{dt}=0$

\paragraph{זמן מחזור בתנועה במימד אחד}

תנועה עם לגרנד'יאן $\mathcal{L}=\frac{1}{2}a\left(q\right)\dot{q}^{2}-V\left(q\right)$
ובקור' קרטזיות $\mathcal{L}=\frac{1}{2}m\dot{x}^{2}-V\left(x\right)$
אחרי חישוב האנרגיה במערכת, נוכל לקבל זמן של תנועה מ$x_{0}$ ל$x$
ע“י האינטגרל:
\[
t=\int_{x_{0}}^{x}\frac{dx}{\sqrt{\frac{2}{m}\left(E-V\left(x\right)\right)}}
\]
ובהתאם, את זמן המחזור בין שתי נקודות שוות אנרגיה פוטנציאלית:
\[
\tau=2\int_{x_{1}\left(E\right)}^{x_{2}\left(E\right)}\frac{dx}{\sqrt{\frac{2}{m}\left(E-V\left(x\right)\right)}}
\]


\paragraph{תנע מוכלל}

$p_{i}=\frac{\partial\mathcal{L}}{\partial\dot{q}_{i}}$. אם קורדינטה
היא \textbf{ציקלית} )כלומר $\frac{\partial\mathcal{L}}{\partial q_{i}}=0$(
נקבל \textbf{שימור בתנע })$p_{i}=\text{const}$(. במקרה זה, נוכל לקבל
לגרנג'יאן חדש, אשר לא תלוי ב$q_{i}$ אלה רק ב$p_{i}$. נעשה זאת באמצעות
טרנספורם לג'נדר )עבור $k$ קור' ציקליות(: $L^{\ast}=L-\sum_{i=1}^{k}p_{i}\dot{q}_{i}$

\paragraph{משפט נתר:}

אם הלגרנז'יאן אינווריאנטי תחת טרנספורמציה רציפה של הקור', כלומר:
\[
L_{s}\equiv L\left(Q\left(s,t\right),\dot{Q}\left(s,t\right),t\right)=L\left(q\left(t\right),\dot{q}\left(t\right),t\right)
\]
קיים גודל שמור עבור כל פרמטר רציף $s_{j}$ של הטרנספורמציה:
\[
I_{j}\left(q_{1},...,q_{N},\dot{q}_{1}...\dot{q}_{N}\right)=\sum_{k=1}^{N}p_{k}\left.\frac{\partial Q_{k}}{\partial s_{j}}\right|_{s_{j}=0}=\unit{const}
\]
\textbf{מפשט נתר המוכלל:} אם $L_{s}=L\left(q\left(t\right)\dot{q}\left(t\right),t\right)+\frac{dF\left(s\right)}{dt}$
אז הגודל השמור הוא 
\[
I=\sum_{k=1}^{N}p_{k}\left.\frac{\partial Q_{k}}{\partial s}\right|_{s=0}-\left.\frac{dF}{ds}\right|_{s=0}
\]


\paragraph{לגרנגיאנים שחוזרים על עצמם:}
\begin{center}
\begin{tabular}{|c|c|c|}
\hline 
 & $K$ & $V$\tabularnewline
\hline 
\hline 
קור' פולאריות & $\frac{m}{2}\left(\dot{r}^{2}+r^{2}\dot{\phi}^{2}\right)$ & \tabularnewline
\hline 
קור' כדוריות & $\frac{m}{2}\left(\dot{r}^{2}+r^{2}\dot{\theta}^{2}+r^{2}\sin^{2}\theta\dot{\varphi}^{2}\right)$ & \tabularnewline
\hline 
 &  & \tabularnewline
\hline 
 &  & \tabularnewline
\hline 
\end{tabular}
\par\end{center}

\subsection{שיט שרבקה הוסיפה}

\textbf{סקיילינג }שיטה להערכת פתרון המערכת עד כדי קבוע. נשתמש בתכונות
של פונקציות הומגניות.\textbf{ הסבר באמצעות דוגמה}, אם יש לנו אנרגיה
קינטית )הומוגנית מדרגה {\LR{2}}( ופוטנציאל $gr^{k}$ )$k$ הומוגני(,
נוכל לבדוק מה קורה ב“מתיחה“ של הפרמטרים בפונקציה - כלומר:
\[
x\rightarrow\alpha x,\ \ t\rightarrow\beta t\Rightarrow v\rightarrow\frac{\alpha}{\beta}v
\]
נקבל לגרנג'יאן מתוח:
\[
\mathcal{L}=K\left(\vec{v}\right)-V\left(\vec{x}\right)\rightarrow\tilde{\mathcal{L}}=\left(\frac{\alpha}{\beta}\right)^{2}K\left(\vec{v}\right)-\alpha^{k}V\left(\vec{x}\right)
\]
משוואות התנועה נשארות זהות אם מכפילים את הלגרנג'יאן בקבוע, לכן נדרוש
כי 
\[
\left(\frac{\alpha}{\beta}\right)^{2}=\alpha^{k}\Rightarrow\beta=\alpha^{1-\frac{1}{2}k}
\]
וקיבלנו יחס דמיון בין אורכים וזמנים. כלומר למדנו משהו על מערכת: אם
נאריך את הזמן פי $\beta$, נקבל מרחק גדול פי $\alpha^{1-\frac{1}{2}k}$.

\textbf{גדלים נוספים שניתן לפתח בצורה דומה:}\\
\textbf{זמן }$\tau\sim d^{1-\nicefrac{k}{2}}$\textbf{ מהירות $v\sim\frac{d}{\tau}\sim d^{\nicefrac{k}{2}}$}\\
\textbf{אנרגיה $E\sim mv^{2}\sim v^{2}\sim d^{k}$ תנ“ז $L\sim\frac{d^{2}}{\tau}\sim d^{1+\nicefrac{k}{2}}$}

\paragraph{משפט הפאי של בקינגהם}

במערכת פיזיקלית המכילה $n$ גדלים עם יחידות, ו-$r$ יחידות בסיסיות
שמקיימת משוואה מהצורה $f\left(x_{1},...,x_{n}\right)=0$ ישנם $n-r$
גדלים חסרי יחידות $\Pi_{i}$ שאיתם אפשר לרשום את כל המידע הכלול ב$f$
ע“י פונקציה חדשה שמקיימת $\tilde{f}\left(\Pi_{1},...,\Pi_{n-r}\right)$.

\paragraph{מציאת הפוטנציאל במסלול}

בהנתן מסלול $r\left(\theta\right)$, נמצע את הפוטנציאל המרכזי משיקולי
אנרגיה ותנע. ראינו, כי האנרגיה הכוללת בפוטנציאל מרכזי תהיה $E=\frac{m}{2}\dot{r}^{2}+V\left(r\right)+\frac{p_{\theta}^{2}}{2mr^{2}}$.
מכאן ניתן לבודד את הפוטנציאל עם טריקים מתמטיים.

\paragraph{המשפט הויריאלי}

צורה כללית: $2\left\langle K\right\rangle =\left\langle \sum_{j}\frac{\partial\mathcal{L}}{\partial q_{j}}q_{j}\right\rangle $
ועבור תנועה בפוטנציאל $V\left(\vec{r}\right)$ $k$-הומוגני עם תנועה
חסומה, מתקיים $2\left\langle K\right\rangle =k\left\langle V\right\rangle $.

\subsection{הבעיה הדו-גופית}

פתרון הבעיה מתחיל מרידוקציה של הבעיה התלת מימדית לדו-מימדית )שימור
תנ“ז( ולאחר מכן לבעיה חד מימדית )פתרון רק עבור הרדיוס כתלות בזמן(.
מתקבל \textbf{לגרנג'יאן אפקטיבי: }$\mathcal{L}_{eff}=\frac{\mu}{2}\left(\dot{r}+r^{2}\dot{\theta}^{2}\right)-V\left(r\right)$.\\
לעתים נוח לעבור למתיאור המסלול לפי הזמן $r\left(t\right)$ לתיאור
המסלול לפי הזווית $r\left(\theta\right)$. הקשר המתקבל הוא: $\dot{r}=\frac{p_{\theta}}{\mu r^{2}}r^{\prime}$.
נציב את הלגר'נג'יאן בא.ל, נסמן $F\left(r\right)=-\frac{dV}{dr}$,
נחליף משתנה $u=\frac{1}{r}$ ומקבלים את\\
 \textbf{משוואת בינה:} $u^{\prime\prime}+u=-\frac{\mu}{p_{\theta}^{2}u^{2}}F\left(\nicefrac{1}{u}\right)$

\paragraph{בעיית קפלר}

בעיה העוסקת בפוטנציאלים מהצורה $V\left(r\right)=\frac{-\alpha}{r}$,
כאשר למשל בכבידה $\alpha=Gm_{1}m_{2}$. עם פוטנציאל מסוג זה, נקבל
משוואת בינה $u^{\prime\prime}+u=\frac{\mu\alpha}{p_{\theta}^{2}}$
והפתרון הכללי למשוואה הוא 
\[
\frac{1}{r\left(\theta\right)}=u\left(\theta\right)=\frac{\mu\alpha}{p_{\theta}^{2}}+A\cos\left(\theta-\theta_{0}\right)=
\]
נגזיר קבועים $r_{0}$ ו$\epsilon$:
\[
r_{0}=\frac{p_{\theta}^{2}}{\alpha\mu}=\frac{\mu\alpha}{p_{\theta}^{2}},\ \ \epsilon=\sqrt{1+\frac{2\mu E}{p_{\theta}}r_{0}^{2}}=\sqrt{1+\frac{2Ep_{\theta}^{2}}{\mu\alpha^{2}}}
\]
\[
\Rightarrow r\left(\theta\right)=\frac{r_{0}}{1+\epsilon\cos\left(\theta-\theta_{0}\right)}=\frac{a\left(1-\epsilon^{2}\right)}{1+\epsilon\cos\left(\theta-\theta_{0}\right)}
\]

\begin{tabular}{|c|c|c|c|c|}
\hline 
$r_{max}$ & $r_{min}$ & אנרגיה & אקסנטריות & צורה\tabularnewline
\hline 
\hline 
$\unit{const}$ & $\unit{const}$ & -$\frac{-\mu\alpha^{2}}{2L^{2}}$ & $\epsilon=0$ & מעגל\tabularnewline
\hline 
$\frac{r_{0}}{1-\epsilon}$ & $\frac{r_{0}}{1+\epsilon}$ & $E<0$ & $\epsilon<1$ & אליפסה\tabularnewline
\hline 
$\infty$ & $\frac{1}{2}r_{0}$ & $E=0$ & $\epsilon=1$ & פרבולה\tabularnewline
\hline 
$\infty$  & $\frac{r_{0}}{1+\epsilon}$ & $E>0$ & $\epsilon>1$ & היפרבולה\tabularnewline
\hline 
\end{tabular}

\paragraph{משפט ברטרנד}

אומר שקיימים בדיוק שני סוגי פוטנציאלים עבורם כל מסלול חסום הוא גם
סגור - פוטנציאל קפלרי )$\frac{-\alpha}{r}$( ופוט' הרמוני )$\frac{1}{2}kr^{2}$(\\
\textbf{וקטור לפרלס רונגה לנץ} גודל שמור של המערכת, מצביע לאורכך הציר
הראשי של האליפסה וגודלו קבוע לאורך כל התנועה: $\vec{A}=\vec{v}\times\vec{l}-\alpha\hat{r}$\textbf{}\\
\textbf{תנאי לתנועה סגורה על עצמה} מסלול יהיה סגור על עצמו אם מתקיים:
$\theta\left(r_{max}\right)-\theta\left(r_{min}\right)=Q\pi$ כאשר
$Q$ מספר רציונאלי.

\paragraph{זווית פיזור:}

\[
\Delta\theta=2\int_{r_{min}}^{\infty}\frac{\nicefrac{l}{r^{2}}}{\sqrt{2\mu\left(E-V\left(r\right)-\nicefrac{l^{2}}{2\mu r^{2}}\right)}}dr
\]


\subsection{תנודות קטנות}

\textbf{ע.מ למצוא נקודת ש.מ} סביבה נפתח את התנודות הקטנות, נגזור את
הפוט' ונמצא אילו ערכים מאפסים את הנגזרת השניה שלו. יש דרכים מתמטיות
למצוא יציבות, אבל אין מקום לכל הציורים של רות לורנץ אז הכי טוב וקל
לקבוע משיקולים פיזיקאליים.

\paragraph{סכמת פתרון}

מניחים שלמערכת יש ש“מ ב$\vec{q_{0}}$. מגידירים מערכת קור' חדשות
$q_{i}$, לפי הזזה מש.מ. מפתחים את הלגרנג'יאן מסדר שני. \textbf{מזניחים}:
נגזרות מסדר ראשון בזמן או בקור'.קבועים.\\
מבחיאים לצורה:
\[
\mathcal{L}_{eff}\approx\frac{1}{2}\dot{\vec{q}}^{T}M\dot{\vec{q}}-\frac{1}{2}\vec{q}^{T}K\vec{q}
\]
כעת, נפתור את בעיית הערכים העצמיים:
\[
\det\left(K-\omega^{2}M\right)=0
\]
ונקבל פתרונות. הפתרון הסופי יפרש ע“י $\vec{\omega}$ וקטור הע“ע
ו$\vec{u_{i}}$ הוקטורים העצמיים המתאימים: 
\[
\vec{q}\left(t\right)\approx\sum_{i}\vec{u_{i}}\cos\omega_{i}t
\]
\textbf{שיטה מקבילה למצוא את $M$ ו$K$:}\\
\textbf{$K_{ij}=\text{ׂׂ}\left.\frac{\partial^{2}V}{\partial q_{j}\partial q_{i}}\right|_{\vec{q}=0},\ M=\left.\frac{\partial T}{\partial\dot{q}_{j}\partial\dot{q_{i}}}\right|_{\vec{q}=0}$}

\subsection{פונקצית ריילי וחיכוך}

פונקציית ריילי מוגדרת עבור מערכות עם חיכוך לינארי $c_{ij}$ למהירות
בצורה:
\[
\mathcal{R}=\frac{1}{2}\sum_{ij}c_{ij}\dot{q}_{i}\dot{q_{j}}
\]
מכאן, ועם כוחות חיצוניים על המערכת $F$, נקבל משוואת א.ל:
\[
\frac{d}{dt}\left(\frac{\partial\mathcal{L}}{\partial\dot{q}_{i}}\right)-\frac{\partial\mathcal{L}}{\partial q_{i}}+\frac{\partial\mathcal{R}}{\partial\dot{q_{i}}}=F
\]


\subsection{תהודה פרמטרית}

מערכת המקמיית את משוואת התנועה $\ddot{x}+\omega\left(t\right)x=0$
כאשר $\omega$ פונק' מחזורית עם זמן מחזור $T$.

\paragraph{תנאים לתהודה פרמטרית}

נלכסן את מרחב הפתרונות של המשוואה, ונקבל זוג פתרונות. ממחזוריות המשוואה,
נקבל שפתרון $x_{i}$ מקיים: $x_{i}\left(t+nT\right)=\mu_{i}^{n}x_{i}$
כאשר הפתרון למוואה כזו הוא $x_{i}=\left(t\right)=\mu^{\frac{t}{T}}\Pi_{i}\left(t\right)$.
)$\Pi_{i}$ פונקציה מחזורית(. מכאן נקבל שמורה של המערכת - עבור זוג
פתרונות $x_{2}\dot{x_{1}}-x_{1}\dot{x_{2}}=\text{const}$. מכאן נובע
ש$\mu_{1}\mu_{2}=1$. לכן, מתקיימת תהודה פרמטרית אם $\mu_{1}\neq1$.

\subsubsection{רוחב התהודה עבור משוואת \textenglish[variant=american]{mathiu}}

כלומר, מתי עבור תדירות מאולצת $\omega\left(t\right)=\omega_{0}^{2}\left(1+\varepsilon\cos\gamma t\right)$
נקבל תהודה? עבור איזה טווח סביב הרזוננס של המערכת $\gamma=\omega_{0}+\delta$?
צריך להציב במדר, לנחש פתרון מהצורה $a\left(t\right)\cos\left(\omega_{0}+\frac{\delta}{2}\right)t+b\left(t\right)\sin\left(\omega_{0}+\frac{\delta}{2}\right)t$,
מניחים ש$a,b$ קבועים ביחס לתנודות המהירות, מזניחים איברים עם תדירויות
גבוהות ומקבלים תנאי: $\left|\delta\right|<\frac{1}{2}\omega_{0}^{2}\varepsilon$.
)מסדר שני נקבל איבר $\frac{\varepsilon^{2}}{32}\omega_{0}^{2}$.

\paragraph{טריק למציאת נק' ש.מ:}

נתחיל מלהביא את הפתרון לצורה $\ddot{\phi}=-\frac{\partial V}{\partial\phi}+f\left(\phi,t\right)$,
כאשר $V$ פוטנציאל הרמוני ו$f$ הכוח המאלץ האפקטיבי על הבעיה. נפצל
את הבעיה לפתרון איטי ומהיר $\phi=\Phi+\eta$, ונקבל פוטנציאל אפקטיבי:
\[
V_{eff}\left(\Phi\right)=V\left(\Phi\right)+\frac{1}{2}\left\langle \dot{\eta}^{2}\right\rangle 
\]
כאשר $\dot{\eta}\left(\Phi,t\right)=\int_{t_{0}}^{t}f\left(t^{\prime},\Phi\right)dt^{\prime}$.
נמצא את נקודות הקיצון של הפוטנציאל האפקטיבי ע“מ למצוא את נק' ש.מ.

\subsection{פיזור}

\textbf{זווית מרחבית - }הכללה של זווית לקור' כדוריות - $d\Omega=\frac{dS}{r^{2}}=\sin\theta d\theta d\varphi$
בבעיות עם \textbf{סימטריה גלילית} כמו שלמדנו, $d\Omega=2\pi\sin\theta d\theta$.

\paragraph{חתך הפעולה}

השטח ממנו חלקיקים באלומת האור מפזרים לזווית המרחבית $d\Omega$. בהנתן
פרמטר פגיעה $\rho$ כפונקציה של זווית הפיזור $\theta$, ניתן לקבל
ביטוי לחתך הפועולה:
\[
\frac{d\sigma}{d\Omega}=\frac{\rho\left(\theta\right)}{\sin\theta}\left|\frac{d\rho}{d\theta}\right|
\]


\paragraph{חתך הפעולה הכולל}

\[
\sigma_{T}=\int\frac{d\sigma}{d\Omega}d\Omega=2\pi\int_{0}^{\pi}\rho\left(\theta\right)\left|\frac{d\rho}{d\theta}\right|d\theta
\]


\subsubsection{בפוטנציאל מרכזי}

נמצא את התנ“ז במרחקים ארוכים לפי $l=\mu\rho v_{\infty}$., ואת
האנרגיה הכוללת בתור $E=\frac{1}{2}\mu v_{\infty}^{2}$. מכאן נקבל
את הקשר: $r_{min}^{2}\left(1-V\left(r_{min}\right)/E\right)=\rho^{2}$
אותו נציב במשוואת הקשר בין $r$ ל$\theta$ בבעיית שתי הגופים.

\subsection{גוף קשיח}

\subsubsection{דו מימד}

$I=\sum_{\alpha}m_{\alpha}r_{\alpha}^{2}\unit{\ or\ }\int r^{2}dm$

$T=T_{cm}+T_{rot}=\frac{Mv_{cm}^{2}}{2}+\frac{I_{cm}\omega^{2}}{2}$

\subsubsection{תלת מימד}

$T=T_{cm}+T_{rot}=\frac{1}{2}Mv_{cm}^{2}+\frac{1}{2}\vec{\omega}^{T}I\vec{\omega}$\\
$I_{ij}=\sum_{\alpha}m_{\alpha}\left(r_{\alpha}^{2}\delta_{ij}-r_{\alpha,i}r_{a,j}\right)$\\
$\unit{or}\int\left(r^{2}\delta_{ij}-r_{i}r_{j}\right)dm$

\paragraph{תנ“ז )גם לדו-מימד(}

$\vec{J}=M\vec{R_{cm}}\times v_{cm}+I_{cm}\vec{\omega}$\\
\textbf{משפט הצירים המאונכים}\ \ \ עבור גוף דו מימדי במישור $xy$:
$I_{xx}+I_{yy}=I_{zz}$\\
\textbf{סיבוב סביב ציר לא ראשי\ \ \ }גוף בעל מומנט אינרציה $I$
מסתובב סביב ציר $\hat{n}$ בקור' הראשיות של הגוף. יקבל מומנט אינרציה:
$I_{\hat{n}}=\hat{n}^{T}I\hat{n}$.
\begin{center}
\begin{tabular}{|c|c|c|}
\hline 
\textbf{מומנט} & \textbf{ציר} & \textbf{גוף}\tabularnewline
\hline 
\hline 
$\frac{Ml^{2}}{12}$ & מרכז & \textbf{מוט}\tabularnewline
\hline 
$\frac{mL^{2}}{3}$ & קצה & \tabularnewline
\hline 
$\frac{mr^{2}}{2}$ & סימטריה & \textbf{גליל מלא}\tabularnewline
\hline 
$\frac{m\left(3r^{2}+h^{2}\right)}{12}$ & מרכז & \tabularnewline
\hline 
$\frac{mr^{2}}{4}$ & הקוטר & \textbf{דיסקית}\tabularnewline
\hline 
$mr^{2}$ & סימטריה & \textbf{חישוק}\tabularnewline
\hline 
$\frac{mr^{2}}{2}$ & הקוטר & \textbf{חישוק}\tabularnewline
\hline 
$\frac{2}{3}mr^{2}$ &  & \textbf{כדור חלול}\tabularnewline
\hline 
$\frac{2}{5}mr^{2}$ &  & \textbf{כדור מלא}\tabularnewline
\hline 
$\frac{M\left(a^{2}+b^{2}\right)}{12}$ & ציר ניצב & \textbf{תיבה\textbackslash פלטה}\tabularnewline
\hline 
$\frac{Mb^{2}}{12}$ & ציר מקביל & \textbf{פלטה}\tabularnewline
\hline 
\end{tabular}
\par\end{center}

\paragraph{משפט שטיינר}

נתון $I_{ij}$ במערכת מ“מ. נתונה מערכת צירים מוזזת ב$\vec{a}$,
כלומר $\vec{x}^{\prime}=\vec{x}+\vec{a}$. במערכת המוזזת 
\[
I_{ij}^{\prime}=I_{ij}+M\left[a^{2}\delta_{ij}-a_{i}a_{j}\right]
\]


\paragraph{מציאת פוטנציאל עבור כוח כללי תלוי מיקום}

דוגמה מהתרגול, היה נתון כח על כל חלקיק $\left(r\right)$$\delta F$.
נמצא פוטנציאל עבור כל חלקיק לפי $\delta V=\int\delta Fdr$. אחרי זה,
נמצא את הפוטנציאל על כל הגוף, שיהיה הפוטנציאל האפקטיבי של הבעיה:
\[
V=\int dV=\iint\delta Fdr
\]


\paragraph{משוואות אוילר}

$\begin{cases}
I_{1}\dot{\omega_{1}}+\left(I_{3}-I_{2}\right)\omega_{2}\omega_{3} & =\tau_{1}\\
I_{2}\dot{\omega_{2}}+\left(I_{1}-I_{3}\right)\omega_{3}\omega_{1} & =\tau_{2}\\
I_{3}\dot{\omega_{3}}+\left(I_{2}-I_{1}\right)\omega_{1}\omega_{2} & =\tau_{3}
\end{cases}$

\paragraph{סביבון סימטרי\protect \\
}

\begin{align*}
\mathcal{L} & =\frac{1}{2}M\dot{\vec{R^{2}}}+\frac{1}{2}I_{1}\left[\sin^{2}\theta\dot{\phi}^{2}+\dot{\theta_{1}}^{2}\right]\\
 & +\frac{1}{2}I_{3}\left[\dot{\psi}^{2}+2\dot{\psi}\dot{\phi}\cos\theta+\cos^{2}\theta\dot{\phi}^{2}\right]-V
\end{align*}

\[
T=\frac{J^{2}}{2}\left(\frac{1}{I_{1}}+\cos^{2}\theta\left(\frac{1}{I_{3}}-\frac{1}{I_{1}}\right)\right)
\]


\paragraph{סביבון סימטרי ללא אילוץ.}

$T,\theta,\omega_{1},\omega_{2},J,\alpha=\unit{const}$. \\
כמו כן: $\omega_{p}=\frac{J}{I_{1}}$, וגם $\omega_{r}=\omega_{3}-\frac{J}{I_{1}}\cos\theta=J_{3}\frac{I_{1}-I_{3}}{I_{1}I_{3}}$

\paragraph{סביבון סימטרי עם כבידה}

נפתח את מומנטי ההתמד במערכת הקודקוד ) משפט שטיינר(. נציב פוטנציאל
כבידתי $mgl\cos\theta$ בלגרנג'יאן. נקבל קור' ציקליות $\phi$ ו$\psi$,
ונסיק שהתנעים הצמודים להם קבועים:
\[
P_{\psi}=I_{3}W_{3}=L_{3}=\unit{const}
\]
\[
P_{\phi}=I_{3}\left(\dot{\psi}+\dot{\phi}\cos\theta\right)\cos\theta+I_{1}\sin^{2}\theta\dot{\phi}=\unit{const}
\]
ואם נבחר את ציר $z$ לכיוון $\vec{L}$, נקבל $P_{\phi}=\left|\vec{L}\right|$.
כעת נותר לבודד את $\dot{\phi}$ ו$\dot{\psi}$:
\[
\dot{\phi}=\frac{p_{\phi}-p_{\psi}\cos\theta}{\left(Md^{2}+I_{1}\right)\sin^{2}\theta}=\omega_{p}
\]
\[
\dot{\psi}=\frac{p_{\psi}}{I_{3}}-\frac{p_{\phi}-p_{\psi}\cos\theta}{\left(Md^{2}+I_{1}\right)\sin^{2}\theta}
\]
כעת, ניתן לכתוב את האנרגיה הכוללת )שהיא כאמור גם כן שמורה של המערכת(,
להעביר אותה ממונחים של מהירות למונחים של תנע,להחסיר כמה איברים ולקבל
קבוע:
\begin{align*}
\tilde{E} & \equiv\frac{1}{2}\underset{M^{\prime}}{\underbrace{\left(Md^{2}+I_{1}\right)}}\dot{\theta}^{2}\\
 & +\underset{V_{eff}}{\underbrace{\frac{\left(p_{\phi}-p_{\psi}\cos\theta\right)^{2}}{2\left(Md^{2}+I_{1}\right)\sin^{2}\theta}-Mgd\left(1-\cos\theta\right)}}
\end{align*}
את הביטוי ניתן לפרש כתנועה במימד אחד של חלקיק אנרגיה $\tilde{E}$,
עם מסה $M^{\prime}$ ופוט' $V_{eff}$. מכאן ניתן לחלץ זמן מחזור לפי
הנוסחא של זמן מחזור לתנועה במימד אחד.

\subsection{משוואות המילטון}

ההמילטוניאן מוגדר כטרנספורם לג'נדר של הלגרנג'יאן, והוא פונקציה של
הקור' $\left\{ q_{i}\right\} $ והתעים הצמודים להם $\left\{ p_{j}\right\} $:

\[
\mathcal{H}=p_{j}\dot{q}_{j}-\mathcal{L}\left(\left\{ q_{i}\right\} ,\left\{ \dot{q}_{i}\right\} ,t\right)
\]
\begin{align*}
\dot{q_{i}} & =\frac{\partial\mathcal{H}}{\partial p_{i}}=\left\{ \mathcal{H},q_{i}\right\} \\
\dot{p_{i}} & =-\frac{\partial\mathcal{H}}{\partial q_{i}}=-\left\{ \mathcal{H},p_{i}\right\} 
\end{align*}

\textbf{זהויות חשובות:}
\[
\frac{d}{dt}f\left(p_{i},q_{i},t\right)=\left\{ \mathcal{H},f\right\} +\frac{\partial f}{\partial t}
\]
\[
\frac{d}{dt}\left\{ f,g\right\} =\frac{\partial\left\{ f,g\right\} }{\partial t}+\left\{ \left\{ f,g\right\} ,\mathcal{H}\right\} 
\]
ומכאן נקבל משהו חשוב - אם $f,g$ שמורות של המערכת, אז גם $\left\{ f,g\right\} $
שמורה של המערכת.

\subsection{טרנספורמציות-קנוניות}

מתקבלות ע“י גזירה של פונקציה יוצרת, שהיא אחת מבין ארבע סוגים. עבור
כל הסוגים מתקיים $\mathcal{H}^{\prime}=\mathcal{H}+\frac{\partial F}{\partial t}$
\begin{center}
\begin{tabular}{|c|c|}
\hline 
משוואות הטרנספורמציה & פונק יוצרת\tabularnewline
\hline 
\hline 
$p=\frac{\partial F_{1}}{\partial q},\ P=-\frac{\partial F_{1}}{\partial Q}$ & $F_{1}\left(q,Q,t\right)$\tabularnewline
\hline 
$p=\frac{\partial F_{2}}{\partial q},\ Q=\frac{\partial F_{2}}{\partial P}$ & $F_{2}\left(q,P,t\right)$\tabularnewline
\hline 
$q=-\frac{\partial F_{3}}{\partial p},\ P=-\frac{\partial F_{3}}{\partial Q}$ & $F_{3}\left(p,Q,t\right)$\tabularnewline
\hline 
$q=-\frac{\partial F_{4}}{\partial p},\ Q=\frac{\partial F_{4}}{\partial P}$ & $F_{4}\left(p,P,t\right)$\tabularnewline
\hline 
\end{tabular}
\par\end{center}

\textbf{איך עוברים בין פונקציות? }נניח שיש לי פונקציה יוצרת $F_{1}$
ואני רוצה למצוא פונקציה יוצרת מקבילה $F_{3}$, כלומר, ארצה להפוך את
$q$ ל$p$. נעשה טרנספורם לג'נדר: $F_{3}=pq-F_{2}$, כאשר אנחנו מציבים
את $q$ במונחים של $p$ ו$Q$, כפי שניתן לפתח ממשוואות הטרנס' של $F_{1}$.

\paragraph{טרס' קאנוניות וסוגרי פואסון}

טרנס' היא קאנונים אם“ם לכל $f,g$ במרחב הפאזה מתקיים:
\[
\left\{ f,g\right\} _{p,q}=\left\{ f,g\right\} _{P,Q}
\]


\paragraph{תנאי שקול:}

$\begin{cases}
\left\{ P_{i},Q_{j}\right\} _{P,Q}=\left\{ P_{i},Q_{j}\right\} _{p,q}=\delta_{ij}\\
\left\{ P_{i},P_{j}\right\} _{P,Q}=\left\{ P_{i},P_{j}\right\} _{p,q}=0\\
\left\{ Q_{i},Q_{j}\right\} _{P,Q}=\left\{ Q_{i},Q_{j}\right\} _{p,q}=0
\end{cases}$

\subsubsection{שינויים קטנים בפונקציות}

אם $G$ פונקציה יוצרת כלשהי, שינוי קטן בפונקציה $u$ יראה:

\[
\delta u\left(q,p\right)=\varepsilon\left\{ u,G\right\} 
\]
ומכאן, מתקבל הקשר בין שינוי בפונקציה לבין ההמילטוניאן:
\[
\delta u\left(u,p\right)=\varepsilon\left\{ u,H\right\} =\varepsilon\frac{du}{dt}
\]
כלומר המליטוניאן יוצר העתקות בזמן. \textbf{מסקנה נוספת} היא עבור שמורות
במערכת. נניח כי המערכת מקיימת שמורה $A$ שאינה תלויה בזמן, נקבל את
הקשר:
\[
\delta H=\varepsilon\left\{ H,A\right\} =0
\]
כלומר שמורה יוצרת טרנס' קאנונית שהיא טרנספורמצית סימטריה של ההמולטיניאן.

\subsection*{תורת המילטון יעקובי}

שיטה למציאת פונקציה יוצרת כך שההמילטוניאן החדש מתאפס:
\[
\tilde{\mathcal{H}}=\mathcal{H}\left(q_{j},p_{j}t\right)+\frac{\partial F_{2}}{\partial t}\stackrel{!}{=}0
\]
מכאן נקבל קורדינטות פשוטות מהן, לפי משוואות המילטון $Q_{j}=\alpha_{j}=\unit{const}$
וגם $P_{j}=\beta_{j}=\unit{const}$. \textbf{פתרון כללי} יננתן ע“י
פונקציה יוצרת מהצורה $F_{2}=S\left(q_{1},q_{2},...,q_{n},\alpha_{1},...,\alpha_{n},t\right)$.הפונקציה
\textenglish[variant=american]{\textbf{S}}\textbf{ היא הפעולה:} $S=\int_{t_{1}}^{t}\mathcal{L}\left(q,\dot{q},t\right)\unit{d}t$.
מהגדרדתה פונקציה יוצרת, נקבל את \textbf{משוואות המילטון יעקובי}:
\[
\mathcal{H}\left(q_{j},\frac{\partial S}{\partial q_{j}},t\right)+\frac{\partial S}{\partial t}\stackrel{!}{=}0
\]

\textbf{{\LR{1}}}\textbf{. פותרים את} משוואת המילטון יעקובי ע“י
הפרדת משתנים )בחיבור!(, ומקבלים פתרון $S\left(q,t;\alpha\right)$.
כאשר $\alpha$ אוסף קבועים כלהם של המערכת.\\
\textbf{{\LR{2}}}\textbf{. מוצאים את הקבועים} $\beta$ לפי $\frac{\partial S}{\partial\alpha}=\beta$
)ראינו שנגזרת של $S$ בכל קבוע תנועה של המערכת ייתן קבוע תנועה חדש(,
ומוצאים ביטוי ל$q_{i}\left(\alpha,\beta,t\right)$\\
\textbf{{\LR{3}}}\textbf{. פותרים את} $p=\frac{\partial S}{\partial q}$
ומוצאים את $p_{i}$ - כעת פתרנו את התרגיל עד כדי קבועים. \\
\textbf{{\LR{4}}}\textbf{.} משמשים ב\textbf{תנאי התחלה} כדי למצוא
את $\alpha$ ו$\beta$.\\


\paragraph{טיפים}

אם $\mathcal{H}$ לא תלוי בזמן, נוכל לנחש $S\left(q_{j},\frac{\partial S}{\partial q_{j}},t\right)=W\left(q_{j},\frac{\partial S}{\partial q_{j}}\right)-Et$

\end{multicols}
\end{document}
